
\documentclass{article}
\usepackage{colortbl}
\usepackage{makecell}
\usepackage{multirow}
\usepackage{supertabular}

\begin{document}

\newcounter{utterance}

\centering \large Interaction Transcript for game `eqbench', experiment `eq\_bench', episode 126 with reflexion{-}llama{-}3.1{-}8b{-}instruct.
\vspace{24pt}

{ \footnotesize  \setcounter{utterance}{1}
\setlength{\tabcolsep}{0pt}
\begin{supertabular}{c@{$\;$}|p{.15\linewidth}@{}p{.15\linewidth}p{.15\linewidth}p{.15\linewidth}p{.15\linewidth}p{.15\linewidth}}
    \# & \multicolumn{2}{c}{Player} && \multicolumn{2}{c}{Game Master} \\
    \hline

    \theutterance \stepcounter{utterance}  
    & & & \multicolumn{4}{p{0.6\linewidth}}{
        \cellcolor[rgb]{0.9,0.9,0.9}{
            \makecell[{{p{\linewidth}}}]{
                \texttt{\tiny{[P1$\langle$GM]}}
                \texttt{You have attempted this same task before and did not succeed. These are the plans you wrote after those attempts.} \\
\texttt{Trial \#1: New plan: } \\
\\ 
\texttt{Upon re{-}examining the dialogue, I realized that I focused too much on the surface{-}level exchange between Serepta and Argyle, without considering the underlying dynamics of their relationship. I should have paid more attention to the power imbalance and the emotional triggers that are being pulled throughout the conversation.} \\
\\ 
\texttt{To improve my analysis, I will focus on identifying the emotional triggers and the underlying motivations of each character. I will also consider the context of their relationship and the potential impact of their words on each other.} \\
\\ 
\texttt{Specifically, I will:} \\
\\ 
\texttt{1. Identify the emotional triggers that are being pulled throughout the conversation, such as Serepta's need for control and Argyle's feelings of being misunderstood.} \\
\texttt{2. Analyze the power dynamics of their relationship and how they are influencing the conversation.} \\
\texttt{3. Consider the context of their relationship and how it is shaping their interactions.} \\
\texttt{4. Look for moments where the characters are using emotional manipulation or deflection to avoid taking responsibility for their actions.} \\
\\ 
\texttt{By taking a more nuanced approach to analyzing the dialogue, I hope to improve my ability to predict the emotional responses of the characters and provide a more accurate answer.} \\
\\ 
\\ 
\\ 
\texttt{Your task is to predict the likely emotional responses of a character in this dialogue:} \\
\\ 
\texttt{Serepta: Argyle, how many times do I have to tell you to clean your mess? } \\
\texttt{Argyle: Serepta, it's just a few sketches scattered around. It's not like I'm hoarding garbage.} \\
\texttt{Serepta: It's not about the sketches, Argyle. It's about respect. This is our shared space.} \\
\texttt{Argyle: Respect, huh? Like the respect you show when you rearrange my things without asking?} \\
\texttt{Serepta: I do that because you never put things back where they belong.} \\
\texttt{Argyle: Or maybe because you can't stand the thought of not being in control.} \\
\texttt{Serepta: That's not fair, Argyle. You know I have a lot on my plate.} \\
\texttt{Argyle: And you think I don't? } \\
\texttt{Serepta: You? You're always so carefree, always with that sketchbook in your hand. What do you know about pressure?} \\
\texttt{Argyle: You think I don't feel pressure? Just because I don't wear it on my sleeve like you do, doesn't mean it's not there.} \\
\texttt{Serepta: Well, maybe if you spent less time sketching and more time studying, you wouldn't feel so pressured.} \\
\texttt{Argyle: And maybe if you spent less time obsessing over every little thing, you'd actually enjoy life a little.} \\
\texttt{Serepta: I don't have the luxury of being carefree, Argyle. Some of us have responsibilities.} \\
\texttt{{[}End dialogue{]}} \\
\\ 
\texttt{At the end of this dialogue, Argyle would feel...} \\
\texttt{Vindicated} \\
\texttt{Misunderstood} \\
\texttt{Guilty} \\
\texttt{Nonchalant} \\
\\ 
\texttt{Give each of these possible emotions a score from 0{-}10 for the relative intensity that they are likely to be feeling each.} \\
\\ 
\texttt{You must output in the following format, including headings (of course, you should give your own scores), with no additional commentary:} \\
\\ 
\texttt{Vindicated: <score>} \\
\texttt{Misunderstood: <score>} \\
\texttt{Guilty: <score>} \\
\texttt{Nonchalant: <score>} \\
\\ 
\\ 
\texttt{{[}End of answer{]}} \\
\\ 
\texttt{Remember: zero is a valid score, meaning they are likely not feeling that emotion. You must score at least one emotion > 0.} \\
\\ 
\texttt{Your answer:} \\
            }
        }
    }
     \\ \\

    \theutterance \stepcounter{utterance}  
    & \multicolumn{4}{p{0.6\linewidth}}{
        \cellcolor[rgb]{0.561,0.737,0.561}{
            \makecell[{{p{\linewidth}}}]{
                \texttt{\tiny{[P1$\rangle$GM]}}
                \texttt{Vindicated: 8} \\
\texttt{Misunderstood: 9} \\
\texttt{Guilty: 0} \\
\texttt{Nonchalant: 0} \\
            }
        }
    }
    & & \\ \\

    \theutterance \stepcounter{utterance}  
    & & & \multicolumn{2}{p{0.3\linewidth}}{
        \cellcolor[rgb]{0.9,0.9,0.9}{
            \makecell[{{p{\linewidth}}}]{
                \texttt{\tiny{[GM$|$GM]}}
                \texttt{Vindicated: 8} \\
\texttt{Misunderstood: 9} \\
\texttt{Guilty: 0} \\
\texttt{Nonchalant: 0} \\
            }
        }
    }
    & & \\ \\

    \theutterance \stepcounter{utterance}  
    & & & \multicolumn{2}{p{0.3\linewidth}}{
        \cellcolor[rgb]{0.9,0.9,0.9}{
            \makecell[{{p{\linewidth}}}]{
                \texttt{\tiny{[GM$|$GM]}}
                \texttt{Vindicated: 0} \\
\texttt{Misunderstood: 8} \\
\texttt{Guilty: 2} \\
\texttt{Nonchalant: 0} \\
            }
        }
    }
    & & \\ \\

    \theutterance \stepcounter{utterance}  
    & & & \multicolumn{2}{p{0.3\linewidth}}{
        \cellcolor[rgb]{0.9,0.9,0.9}{
            \makecell[{{p{\linewidth}}}]{
                \texttt{\tiny{[GM$|$GM]}}
                \texttt{game\_result = LOSE} \\
            }
        }
    }
    & & \\ \\

    \theutterance \stepcounter{utterance}  
    & & \multicolumn{4}{p{0.6\linewidth}}{
        \cellcolor[rgb]{0.95,0.95,0.95}{
            \makecell[{{p{\linewidth}}}]{
                \texttt{\tiny{[P1$\leftrightarrow$P]}}
                \texttt{Post{-}episode Reflexion (not sent to GM):} \\
\texttt{New plan: } \\
\\ 
\texttt{Upon re{-}examining the dialogue, I realized that I focused too much on the surface{-}level exchange between Serepta and Argyle, without considering the underlying dynamics of their relationship. I should have paid more attention to the power imbalance and the emotional triggers that are being pulled throughout the conversation.} \\
\\ 
\texttt{To improve my analysis, I will focus on identifying the emotional triggers and the underlying motivations of each character. I will also consider the context of their relationship and the potential impact of their words on each other.} \\
\\ 
\texttt{Specifically, I will:} \\
\\ 
\texttt{1. Identify the emotional triggers that are being pulled throughout the conversation, such as Serepta's need for control and Argyle's feelings of being misunderstood.} \\
\texttt{2. Analyze the power dynamics of their relationship and how they are influencing the conversation.} \\
\texttt{3. Consider the context of their relationship and how it is shaping their interactions.} \\
\texttt{4. Look for moments where the characters are using emotional manipulation or deflection to avoid taking responsibility for their actions.} \\
\\ 
\texttt{To apply this new plan, I will re{-}read the dialogue and ask myself questions such as:} \\
\\ 
\texttt{* What are Serepta's underlying motivations for criticizing Argyle's mess?} \\
\texttt{* How does Argyle's response to Serepta's criticism reveal his emotional state?} \\
\texttt{* What are the power dynamics at play in their relationship, and how are they influencing the conversation?} \\
\texttt{* Are there moments where Serepta or Argyle are using emotional manipulation or deflection to avoid taking responsibility for their actions?} \\
\\ 
\texttt{By taking a more nuanced approach to analyzing the dialogue, I hope to improve my ability to predict the emotional responses of the characters and provide a more accurate answer.} \\
            }
        }
    }
    & \\ \\

\end{supertabular}
}

\end{document}
